\section{Podsumowanie}

W ramach projektu aplikacji PocztaExpress zrealizowano kluczowe funkcjonalności umożliwiające użytkownikom wygodne nadawanie przesyłek bez konieczności wizyty w placówkach pocztowych. Aplikacja obsługuje rejestrację użytkowników, historię wysyłek oraz podstawowe operacje logistyczne związane z nadawaniem paczek. System został zaprojektowany z myślą o intuicyjnym interfejsie użytkownika oraz kompatybilności z systemami Windows 10/11 i .NET 6.0.
\newline

\noindent Zrealizowane prace obejmowały:
\begin{itemize}
\item Opracowanie i implementację interfejsu użytkownika (okna logowania, formularz nadania paczki, historia przesyłek itp.).
\item Integrację z bazą danych do przechowywania informacji o przesyłkach i użytkownikach.
\item Implementację funkcji rejestracji, logowania oraz zarządzania kontem użytkownika.
\item Dodanie modułu zmiany hasła i sekcji kontaktowej.
\item Stworzenie systemu wyświetlania historii nadanych paczek.
\item Stworzenie systemu dodawania punktów do użytkownika
\item Możliwość dodania komentarzy
\end{itemize}

\noindent W przyszłości możliwe jest jej rozwinięcie o dodatkowe funkcjonalności, takie jak integracja z systemami ERP dla firm, automatyczne powiadomienia SMS dla odbiorców oraz obsługa międzynarodowych przesyłek, możliwość śledzenia przesyłek w czasie rzeczywistym, optymalizacja wydajności aplikacji oraz usprawnienie obsługi błędów, rozwój wersji mobilnej aplikacji na systemy Android oraz iOS, umożliwienie płatności online za wysłanie przesyłki Dzięki zastosowanym technologiom aplikacja jest nowoczesnym i efektywnym narzędziem do zarządzania procesem wysyłki paczek.

